\documentclass[11pt]{article}
\usepackage[margin=0.92in]{geometry}
\usepackage{amsmath,amssymb,amsthm,mathtools}
\usepackage[T1]{fontenc}
\usepackage{lmodern}
\usepackage{graphicx}
\usepackage{float}
\usepackage{xcolor}
\usepackage[hidelinks]{hyperref}

\newtheorem{theorem}{Theorem}
\newtheorem{lemma}[theorem]{Lemma}
\newtheorem{corollary}[theorem]{Corollary}
\newtheorem{remark}[theorem]{Remark}
\newcommand{\cX}{\mathcal X}
\newcommand{\cH}{\mathcal H}
\newcommand{\Ran}{\operatorname{Ran}}
\newcommand{\clos}[1]{\overline{#1}}
\newcommand{\hot}{\widehat\otimes_2}
\newcommand{\HT}{\mathop{\widehat\bigotimes}\limits}

\title{Sobolev Tensors Always Belong to the Tensor Product\\
of Their Minimal Subspaces}
\author{Candidate full resolution for arXiv:1809.11001}
\date{13 August 2026}

\begin{document}
\maketitle

\begin{center}
\fcolorbox{black}{gray!8}{\parbox{0.88\linewidth}{
\textbf{Status: candidate full solution, likely valid pending human review.}
The conclusion of Proposition 3.7 of the source holds without Assumption 3.5;
therefore the counterexamples conjectured in Remark 3.6 do not exist.}}
\end{center}

\section{The source conjecture and the result}

Let $\Omega=\prod_{j=1}^d\Omega_j$.  Ali and Nouy identify
\[
 H^1(\Omega)=\bigcap_{k=1}^d
 \left(H^1(\Omega_k)\hot\HT_{j\ne k}L^2(\Omega_j)\right),
 \tag{1}
\]
with an equivalent intersection norm.  For $u\in H^1(\Omega)$, they define
$U_j^{\min}(u)$ as the $H^1(\Omega_j)$-closure of all contractions of $u$
against $L^2$ functionals in the other variables.  Proposition 3.7 proves
\[
 u\in\clos{\bigotimes_{j=1}^d{}^{a}U_j^{\min}(u)}^{\,H^1(\Omega)}
 \tag{2}
\]
under Assumption 3.5: the $H^1$-orthogonal projection onto every
$U_j^{\min}(u)$ must be bounded in $L^2$.  Remark 3.6 conjectures that (2)
can fail when this assumption fails \cite{AN1}.

\begin{figure}[H]
  \centering
  \includegraphics[width=\textwidth]{figures/open_problem_crop.png}
  \caption{The definition, Assumption 3.5, conjecture, and conditional
  proposition in arXiv:1809.11001, PDF page 9.}
\end{figure}

The conjecture is false: (2) always holds.  The source projection need not
extend to $L^2$.  Instead, one spectrally truncates the inclusion of the
minimal subspace into its $L^2$ closure.  The truncated projections are
simultaneously contractive in both norms, and their ranges have equivalent
$H^1$ and $L^2$ norms.  This makes finite tensors dense after truncation;
letting the spectral threshold tend to zero proves (2).

We first prove the mechanism for arbitrary Hilbert pairs.  This also makes
clear that compactness of the Sobolev embedding is unnecessary.

\section{A simultaneous spectral-cutoff lemma}

\begin{lemma}[two-norm cutoffs]\label{lem:cutoff}
Let $U$ and $M$ be Hilbert spaces and let $J:U\to M$ be continuous,
injective, and have dense range.  There are closed subspaces
$S_\varepsilon\subset U$ and operators $E_\varepsilon:U\to U$,
$\widetilde E_\varepsilon:M\to M$ for $\varepsilon>0$ such that:
\begin{enumerate}
\item $E_\varepsilon$ and $\widetilde E_\varepsilon$ are orthogonal
projections in $U$ and $M$, respectively, and
$\widetilde E_\varepsilon J=JE_\varepsilon$;
\item both projections are contractions in their own Hilbert norms;
\item on $S_\varepsilon=\Ran E_\varepsilon$,
\[
 \sqrt\varepsilon\,\|x\|_U\leq\|Jx\|_M
 \leq\|J\|\,\|x\|_U;
 \tag{3}
\]
\item $E_\varepsilon\to I_U$ strongly in $U$, and
$\widetilde E_\varepsilon\to I_M$ strongly in $M$, as
$\varepsilon\downarrow0$.
\end{enumerate}
\end{lemma}

\begin{proof}
Put $K=J^*J$.  This is a bounded positive self-adjoint operator on $U$ and
$\ker K=\ker J=\{0\}$.  By the spectral theorem, set
\[
 E_\varepsilon=\mathbf 1_{[\varepsilon,\infty)}(K),
 \qquad S_\varepsilon=\Ran E_\varepsilon.
 \tag{4}
\]
It is an orthogonal projection in $U$, and for $x\in U$,
\[
 \|JE_\varepsilon x\|_M^2
 =\langle KE_\varepsilon x,E_\varepsilon x\rangle_U
 \leq \langle Kx,x\rangle_U=\|Jx\|_M^2.
 \tag{5}
\]
Furthermore, if $x\in S_\varepsilon$, spectral localization gives
$\|Jx\|_M^2\geq\varepsilon\|x\|_U^2$, proving (3).

The range $J(S_\varepsilon)$ is closed in $M$: an $M$-Cauchy sequence in
that range is $U$-Cauchy by (3), and its $U$-limit remains in the closed
subspace $S_\varepsilon$.  Since $E_\varepsilon$ commutes with $K$,
\[
 \langle JE_\varepsilon x,Jy\rangle_M
 =\langle K E_\varepsilon x,y\rangle_U
 =\langle Jx,J E_\varepsilon y\rangle_M.
 \tag{6}
\]
Thus (5) lets $JE_\varepsilon J^{-1}$ extend from the dense subspace
$J(U)$ to the $M$-orthogonal projection $\widetilde E_\varepsilon$ onto
$J(S_\varepsilon)$.

Because $K$ has no spectral mass at zero,
$E_\varepsilon x\to x$ in $U$ for every $x\in U$.  The continuous inclusion
then gives convergence in $M$ on the dense set $J(U)$; the contraction bound
and density extend it to every element of $M$.
\end{proof}

\section{The abstract intersection theorem}

Let $A_j\hookrightarrow H_j$ be continuous injective dense embeddings of
Hilbert spaces, $1\leq j\leq d$.  All tensor products below are canonical
Hilbert tensor products.  Inside
$\cH=\HT_{j=1}^d H_j$, set
\[
 X_k=A_k\hot\HT_{j\ne k}H_j,
 \qquad
 \cX=\bigcap_{k=1}^dX_k,
 \qquad
 \|v\|_{\cX}=\max_k\|v\|_{X_k}.
 \tag{7}
\]
For $u\in\cX$, let $H_{[j]}=\HT_{i\ne j}H_i$ and define
\[
 U_j=\clos{\{(I_{A_j}\otimes\ell)u:
                \ell\in H_{[j]}^*\}}^{\,A_j}.
 \tag{8}
\]
It is equivalent in (8) to use only finite linear combinations of
decomposable functionals, since these are dense in $H_{[j]}^*$ and the
contraction map $H_{[j]}^*\to A_j$ is continuous.

\begin{theorem}[minimal subspaces in a Hilbert intersection]\label{thm:abstract}
In the setting above, every $u\in\cX$ satisfies
\[
 u\in\clos{\bigotimes_{j=1}^d{}^{a}U_j}^{\,\cX}.
 \tag{9}
\]
\end{theorem}

\begin{proof}
We split the proof into three steps.

\medskip\noindent
\textbf{Step 1: put all passive factors into their $H$-closures.}
Let
\[
 M_j=\clos{U_j}^{\,H_j}.
 \tag{10}
\]
Viewing $u$ in $X_k=A_k\hot H_{[k]}$, the usual Hilbert-tensor minimality
argument gives
\[
 u\in U_k\hot H_{[k]}.
 \tag{11}
\]
For completeness, if $P_k$ is the $A_k$-orthogonal projection onto $U_k$,
then every contraction of $((I-P_k)\otimes I)u$ against an element of
$H_{[k]}^*$ vanishes by (8); hence the tensor itself vanishes, proving (11).

Let $Q_j$ be the $H_j$-orthogonal projection onto $M_j$.  The same
contraction argument in $\cH$ shows $(Q_j\otimes I)u=u$ for every $j$, so
\[
 u\in\HT_{j=1}^dM_j.
 \tag{12}
\]
Fix $k$.  For $j\ne k$, $Q_j$ acts contractively on the corresponding
$H_j$ factor of $U_k\hot H_{[k]}$.  Applying their tensor product to (11),
and using (12), gives
\[
 u\in Y_k:=U_k\hot\HT_{j\ne k}M_j
 \quad (1\leq k\leq d).
 \tag{13}
\]
Equality first holds in the common ambient space $\cH$ and therefore in
$X_k$, whose embedding into $\cH$ is injective.

\medskip\noindent
\textbf{Step 2: construct one cutoff that converges in every directional
norm.}
Apply Lemma~\ref{lem:cutoff} to the inclusion
$J_j:U_j\hookrightarrow M_j$.  Denote the resulting projections by
$E_{j,\varepsilon}$ on $U_j$ and $\widetilde E_{j,\varepsilon}$ on $M_j$,
and their common range by $S_{j,\varepsilon}\subset U_j$.
On $Y_k$ define
\[
 T_{k,\varepsilon}
 =E_{k,\varepsilon}\hot
  \HT_{j\ne k}\widetilde E_{j,\varepsilon}.
 \tag{14}
\]
It is a contraction: the active factor is contractive in $A_k$, and every
passive factor is contractive in $H_j$.  Strong convergence of finite Hilbert
tensor products and (13) yield
\[
 T_{k,\varepsilon}u\longrightarrow u
 \quad\text{in }X_k.
 \tag{15}
\]
The tensors in (14) are all restrictions of the same $H$-orthogonal product
projection $\HT_j\widetilde E_{j,\varepsilon}$ on $\HT_jM_j$.
Consequently their values at $u$ agree in $\cH$; call this common value
$u_\varepsilon$.  Since there are finitely many $k$, (15) gives
\[
 \|u_\varepsilon-u\|_{\cX}\longrightarrow0.
 \tag{16}
\]

\medskip\noindent
\textbf{Step 3: approximate each cutoff by finite tensors.}
For fixed $\varepsilon$, Lemma~\ref{lem:cutoff} says that the $A_j$ and
$H_j$ norms are equivalent on $S_{j,\varepsilon}$.  Hence the $d$ Hilbert
tensor norms
\[
 S_{k,\varepsilon}^{A_k}\hot
 \HT_{j\ne k}S_{j,\varepsilon}^{H_j},
 \qquad 1\leq k\leq d,
 \tag{17}
\]
are mutually equivalent on the algebraic tensor product
$\bigotimes_j^a S_{j,\varepsilon}$.  Their completions therefore coincide,
with equivalent norms.  By (14), $u_\varepsilon$ belongs to that common
completion.  It can thus be approximated simultaneously in every $X_k$ norm
by finite tensors from
\[
 \bigotimes_{j=1}^d{}^aS_{j,\varepsilon}
 \subseteq\bigotimes_{j=1}^d{}^aU_j.
 \tag{18}
\]
Choose $\varepsilon_n\downarrow0$ and, using (18), choose $v_n$ with
$\|v_n-u_{\varepsilon_n}\|_{\cX}<1/n$.  Equation (16) then gives
$v_n\to u$ in $\cX$, which is (9).
\end{proof}

\section{Resolution of the Sobolev conjecture}

\begin{corollary}[unconditional form of Proposition 3.7]\label{cor:sob}
In the setting of arXiv:1809.11001, let
$\Omega=\prod_{j=1}^d\Omega_j$ and $u\in H^1(\Omega)$.  Then
\[
 u\in\clos{\bigotimes_{j=1}^d{}^aU_j^{\min}(u)}^{\,H^1(\Omega)}.
 \tag{19}
\]
No $L^2$-boundedness assumption on the $H^1$-orthogonal projections is
required.
\end{corollary}

\begin{proof}
Take $A_j=H^1(\Omega_j)$ and $H_j=L^2(\Omega_j)$.  The source's
identification (1) is precisely (7), up to an equivalent intersection norm.
The source defines $U_j^{\min}(u)$ using decomposable $L^2$ functionals in
the other factors.  As observed after (8), completing those functionals gives
the same $A_j$-closed subspace.  Theorem~\ref{thm:abstract} now gives (19).
\end{proof}

Thus the membership statement conjectured to fail in Remark 3.6 is in fact
universal.  This does not say that Assumption 3.5 itself is automatic: the
full $H^1$-orthogonal projection onto $U_j^{\min}(u)$ can still be unbounded
in $L^2$.  The proof only constructs an increasing family of two-norm-bounded
spectral cutoffs, which is enough for approximation of the single tensor $u$.

\section{Checks, scope, and novelty}

The proof does not use compactness of $H^1(\Omega_j)\hookrightarrow
L^2(\Omega_j)$, eigenfunction expansions, or a uniform bound as
$\varepsilon\downarrow0$.  Continuous spectrum causes no problem: strong
convergence in Lemma~\ref{lem:cutoff} follows from $\ker(J_j^*J_j)=0$.
The only structural inputs are the finite intersection (1), Hilbert tensor
products, and continuous injective component embeddings.

Part II repeats that membership was known only under conditions and says the
authors did not believe it held generally \cite{AN2}.  The earlier work of
Falc\'o and Hackbusch proves density of the algebraic product of the full
component Sobolev spaces in the Sobolev intersection \cite{FH}.  That density
does not put the factors into the tensor-dependent minimal subspaces; the
spectral cutoff in this packet supplies exactly that missing localization.
A bounded search through 13 August 2026 found the two source papers and the
general minimal-subspace literature, but no later paper giving (19) or the
abstract cutoff theorem above.  This is bounded evidence, not an exhaustive
novelty claim.

Human review should focus on three points: (i) the refinement (13) in the
common $L^2$ ambient tensor product; (ii) the $M$-orthogonal extension in
Lemma~\ref{lem:cutoff}; and (iii) the identification of the completions in
(17).  Each is proved explicitly above.

\begin{thebibliography}{9}

\bibitem{AN1}
M.~Ali and A.~Nouy,
\emph{Singular Value Decomposition in Sobolev Spaces: Part I},
arXiv:1809.11001v3 (2020).

\bibitem{AN2}
M.~Ali and A.~Nouy,
\emph{Singular Value Decomposition in Sobolev Spaces: Part II},
arXiv:1912.11293 (2019).

\bibitem{FH}
A.~Falc\'o and W.~Hackbusch,
\emph{On minimal subspaces in tensor representations},
Foundations of Computational Mathematics \textbf{12} (2012), 765--803.

\end{thebibliography}

\end{document}
